\subsection{State estimation with Kalman filter}

As in the pole-placement controller, only the ball position and the coil current are measured.
The velocity required by the state-feedback law is therefore estimated with a continuous-time Kalman
filter.

The estimator is designed on the linearized mechanical model
\begin{equation}
    \dot{\vect{x}}_{kf}
    =
    A_x \vect{x}_{kf}
    +
    B_x \Delta i
    \qquad
    y = C_{meas}\vect{x}_{kf} + v
\end{equation}
where
\begin{equation}
    \vect{x}_{kf}
    =
    \begin{bmatrix}
        \Delta x & \Delta \dot{x}
    \end{bmatrix}^{T}
    \qquad
    C_{meas}
    =
    \begin{bmatrix}
        1 & 0
    \end{bmatrix}
\end{equation}


The tuning parameters are
\begin{equation}
    Q_{kf} = \sigma_a^2,
    \qquad
    R_{kf} = \sigma_{yx}^2
    \qquad
    \sigma_a = 20~\mathrm{m/s^2}
    \qquad
    \sigma_{yx} = 0.1 \times 10^{-3}~\mathrm{m}
\end{equation}

The covariance $Q_{kf}$ accounts for mechanical model uncertainty, represented as an equivalent
acceleration disturbance, while $R_{kf}$ describes the position-sensor noise. Their tuning sets the
trade-off between estimation speed and noise attenuation.

The continuous-time Kalman gain is computed in MATLAB as
\begin{equation}
    L_{kf}
    =
    \mathrm{lqe}(A_x,G_{kf},C_{meas},Q_{kf},R_{kf}),
\end{equation}
leading to the estimator dynamics
\begin{equation}
    \dot{\hat{\vect{x}}}_{kf}
    =
    \left(A_x - L_{kf}C_{meas}\right)\hat{\vect{x}}_{kf}
    +
    B_x \Delta i
    +
    L_{kf}\Delta x_{meas}.
\end{equation}

The resulting gain and estimator poles are
\begin{equation}
    L_{kf}
    = 10^{5} \cdot
    \begin{bmatrix}
        0.0064  \\
        2.0282 
    \end{bmatrix},
    \qquad
    p_{kf}
    =
    -318.45 \pm j\,314.02.
\end{equation}

